\documentclass[10pt,a4paper]{article}

\usepackage[T1]{fontenc}
\usepackage[utf8]{inputenc}
\usepackage{textcomp}\usepackage{lmodern,amsmath,amssymb,array,tabularx,booktabs,xcolor,enumitem}
\usepackage[margin=20mm]{geometry}
\usepackage[unicode,colorlinks=true,urlcolor=black,linkcolor=black]{hyperref}
\usepackage{xcolor}
\usepackage[normalem]{ulem}
\usepackage{kotex}

\setlength{\parindent}{0pt}
\setlength{\parskip}{5pt}
\setlength{\emergencystretch}{2em}
\setlength{\tabcolsep}{4pt}
\renewcommand{\arraystretch}{1.10}
\newcommand{\hd}[1]{\par\vspace{8pt}\textbf{#1}\par\vspace{3pt}}
\newcommand{\circnum}[1]{\textcircled{\scriptsize #1}}
\newcommand{\note}{\textreferencemark\ }
\newcommand{\red}[1]{\textcolor{red}{#1}}
\newcommand{\under}[2]{\underbrace{#1}_{\text{#2}}}
\newcommand{\paper}[1]{\href{#1}{paper}}
\newcolumntype{Y}{>{\raggedright\arraybackslash}X}

%%%%%%%%%%%%%%%%%%%%%%%%%%%%%%%%%%%%%%%%%%%%%%%%%%%%%%
\begin{document}



%%%%%%%%%%%%%%%%%%%%%%%%%%%%%%%%%%%%%%%%%%%%%%%%%%%%%%
\hd{1. Core model: City selection + vehicle deployment + multiple penalties}

\[
\boxed{\begin{aligned}
\min_{y,q,O,U}\quad
&\under{\sum_c F_c y_c}{Entry costs}
+\under{\sum_c K_c q_c}{Vehicle costs}\\
&+\frac{1}{|S|}\sum_s\sum_c\left[
\under{c_c^+O_c^s}{Remaining capacity}
+\under{c_c^-U_c^s}{Foregone contribution margin}
\right].
\end{aligned}}
\]

%%%%%%%%%%%%%%%%%%%%%%%%%%%%%%%%%%%%%%%%%%%%%%%%%%%%%

In $c^+,c^-$, $+,-$ distinguish excess and shortage, rather than the signs of costs. All demand, throughput, and costs use \textbf{the same evaluation period.} \note I suppose that setting the evaluation period is needed, and the model determines the city/vehicle deployment to apply during that period. 
%%%%%%%%%%%%%%%%%%%%move from pg3 to pg1%%%%%%%%%%%%%
\\

\textbf{2. Constraints }

\[
\boxed{\begin{aligned}
0\le q_c\le\bar q_c y_c&\qquad\forall c,\\
\sum_c q_c\le Q,&\\
\kappa_c q_c-M_c p_c^s=O_c^s-U_c^s
&\qquad\forall c,s,\\
y_c\in\{0,1\},\quad q_c\in\mathbb Z_{\ge0},&\\
O_c^s,U_c^s\ge0.&
\end{aligned}}
\]

{\small
\begin{tabularx}{\linewidth}{@{}>{\raggedright\arraybackslash}p{0.32\linewidth}Y@{}}
\toprule
Constraint & Meaning \\ \midrule
$0\le q_c\le\bar q_c y_c$ & Zero vehicles if the city is not entered; up to the city-specific limit if it is entered \\
$\sum_cq_c\le Q$ & Total vehicle limit \\
$\kappa_c q_c-M_c p_c^s=O_c^s-U_c^s$ & The difference between service capacity and robotaxi demand equals unused capacity minus unmet demand. \\
%%%%%%%%%%%%%%%%%%%%%%%%%%%%%%%%%%%%%%%%%%%%%
%%%%%%%%%%%%%% 삭제 예정%%%%%%%%%%%%%%%%%%%%%%
%%%%%%%%%%%%%%%%%%%%%%%%%%%%%%%%%%%%%%%%%%%%%%

\bottomrule
\end{tabularx}
}


%%%%%%%%%%%%%%%%%%%%%%%%%%%%%%%%%%%%%%%%%%%%%

\hd{3. Explanations}

\textbf{A. Construct two types of demand }

\[
D_c^s = M_c p_c^s,\qquad 
D_c^{N,s}= M_c (1-p_c^s)
\]

{\small
\begin{tabularx}{\linewidth}{@{}>{\raggedright\arraybackslash}p{0.12\linewidth}>{\raggedright\arraybackslash}p{0.42\linewidth}Y@{}}
\toprule
Symbol & Meaning & Explanation \\ \midrule
$M_c$ &Total number of \textbf{potential travel opportunities across all modes} & Robotaxi demand + human-driven taxi demand + other demand \\
$p_c^s$ & Probability of choosing robotaxi service for a travel opportunity
in city $c$ under the specified service conditions,
before applying fleet-capacity constraints & Need to be specified with reference service conditions such as fares and waiting times \\
$D_c^s$ & Robotaxi choice demand & Demand \textbf{before} applying vehicle shortages \\
$D_c^{N,s}$ & Demand not choosing robotaxis & Human-driven taxis, other modes, giving up the trip, etc. \\
\bottomrule
\end{tabularx}
}

\textbf{B. Constraints}\\
1) Once demand $D_c^s$ and service capacity $\kappa_c q_c$ are determined, the following two values are calculated based on the deployment outcome; they are not variables that need to be investigated separately. $(x)^+=\max\{x,0\}$. Both values are measured in trips.

\[
\begin{aligned}
%%%%%%%%%%%%%%%%%%%%%%%%%%%%%%%%%%%%%%
O_c^s&=(\kappa_c q_c-D_c^s)^+&&\text{Remaining service capacity},\\
U_c^s&=(D_c^s-\kappa_c q_c)^+&&\text{Unserved potential demand}.
\end{aligned}
\]\\

%%%%%%%%%%%%%%%%%%%%%%%%%2pg%%%%%%%%%%%%%%%%%%%%%%%%%%%%%
\newpage
\hd{4. Notation used in the model}

{\small
\begin{tabularx}{\linewidth}{@{}>{\raggedright\arraybackslash}p{0.1\linewidth}>{\raggedright\arraybackslash}p{0.37\linewidth}Y@{}}
\toprule
Symbol & Meaning & Type and unit \\ \midrule
$C$ & Set of potential candidate cities & Input set; cities \\
$S$ & Set of demand scenarios & Input set; scenarios \\
$y_c$ & 1 if city $c$ is entered, else 0 & Decision variable; binary (0 or 1) \\
$q_c$ & Number of vehicles deployed in city $c$ & Decision variable; vehicles (integer) \\
$\kappa_c$ & Number of trips one vehicle can serve during the period & Scale parameter; trips/vehicle over the evaluation period.  \\
$Q$ & Total number of vehicles available for deployment & Input parameter; vehicles \\
$\bar q_c$ & Upper bound on vehicle deployment in each city & Input parameter; vehicles \\

%%%%%%%%%%%%%% 비용 항목 %%%%%%%%%%%%%%%%%
$F_c$
& Fixed cost of city entry and basic infrastructure 
& Input parameter; currency over the evaluation period
\\

$K_c$
& Per-vehicle preparation, deployment, and holding cost
over the evaluation period
& Input parameter; currency/vehicle over the evaluation period
\\

$c_c^+$
& Additional loss per unit of unused service capacity
& Input parameter; currency/trip of unused capacity
\\

$c_c^-$
& Foregone contribution margin per unit of unserved
potential demand
& Input parameter; currency/unserved trip
\\

$M_c$ & Total number of \textbf{potential travel opportunities across all modes}
in city $c$ during one year. $c$ & Input parameter; usage opportunities over the evaluation period. \newline data or assumption \\


%%%%%%%%%%%%%%%%%%%%%%%%%%%%%%%%%%%%
$p_c^s$ & Probability of choosing robotaxi service for a travel opportunity
in city $c$ under the specified service conditions,
before applying fleet-capacity constraints. & Scenario input; dimensionless probability in $[0,1]$ \\
$D_c^s$ & Robotaxi choice demand & Derived input; trips over the evaluation period. \newline $D_c^s= M_cp_c^s$ \\
$D_c^{N,s}$ & Demand not choosing robotaxis (human-driven taxis, other modes, giving up the trip, etc.) & Derived input; trips over the evaluation period. \newline $D_c^{N,s}= M_c(1-p_c^s)$ \\

$O_c^s$ & Unused service capacity
& Outcome variable; trips over the evaluation period. \newline $O_c^s =(\kappa_c q_c-D_c^s)^+$ \\

$U_c^s$ & Unserved robotaxi demand
& Outcome variable; trips over the evaluation period. \newline $U_c^s=(D_c^s-\kappa_c q_c)^+$ \\


\bottomrule
\end{tabularx}
}

\note Demand-related variables consist of the inputs $p_c^s,M_c$ and outcome variables calculated from scenario demand and deployment in the model. For the baseline model, $M_c$ is fixed.
\note Outcome variables are not quantities to be predicted by collecting data for each one separately. They are outcomes calculated from demand and deployment.



\end{document}
